Abstractive summarizers write fluent text that can say things the source does not. Human studies of news summarization found unsupported content in a large share of system outputs \citep{maynez2020faithfulness}, and the standard overlap metrics barely register it \citep{kryscinski2019critical,goodrich2019factual}. Automatic consistency metrics were built to close that gap: SummaC aggregates natural language inference scores over sentence pairs \citep{laban2022summac}, and QAFactEval asks questions of the summary and checks the answers against the source \citep{fabbri2022qafacteval}. Both were validated against human judgements collected on short news articles \citep{fabbri2021summeval}, and papers now report them next to, or in place of, human evaluation.

Summarization research has meanwhile moved to long inputs: government reports, books and scientific articles several thousand tokens long \citep{huang2021govreport,kryscinski2022booksum}. A long source changes both sides of the comparison. A summarizer has more material to misattribute, merge or invent, and a metric has to find the supporting evidence for each summary sentence somewhere in a document it may have to split into pieces. Human judgement changes as well: annotators need more time per summary and disagree more when the evidence is spread across many pages \citep{krishna2023longeval}. Whether metrics validated on news still order systems the way annotators do on long documents is therefore an empirical question, and it matters most when the metric is the only evaluation a paper reports. Human evaluation of long summaries is slow and expensive, so the temptation to report a metric alone is strongest exactly where the metric has been checked least.

We study that question on an evaluation we had already run for a different purpose. Three systems summarized 200 long documents: lead-3, an extractive baseline; bart-large, an abstractive model; and bart-large-rl, the same model fine-tuned with a consistency reward. Two annotators judged every summary as consistent or inconsistent with its source, and a third resolved disagreements. We scored the same summaries with SummaC and QAFactEval and compared the two views at two levels. At the system level we ask whether the metrics order the three systems as the human consistency rate does. At the summary level we ask how well each metric separates the summaries the annotators accepted from those they rejected. We repeat both comparisons in three buckets of document length, because a failure that only appears on the longest documents would be hidden in the average. The design is small on purpose: three systems and two metrics are enough to test whether the conclusions people draw from these metrics survive a change of document length, and a binary judgement per summary keeps the human side of the comparison affordable.

The system-level answer is reassuring. Both metrics reproduce the human ordering of the three systems overall and in every length bucket, with the extractive baseline first and the consistency-rewarded model second. The scale is another matter: a metric score and a human consistency rate are different quantities, and the gap between them grows with document length. The summary-level answer is less reassuring. The discrimination of SummaC falls steadily from the short to the long bucket, while QAFactEval changes little, so the metric that is better on short documents is the weaker one on long ones. A paper that uses these metrics to compare systems on long documents can rely on their ranking; one that uses them to flag individual unfaithful summaries should check how long its documents are.

Our contributions are the following.
\begin{itemize}
\item A human consistency annotation of the summaries of 200 long documents produced by three systems, with two annotators per summary, adjudication, and the length of every source recorded.
\item A comparison of SummaC and QAFactEval with the human judgements at the system level and at the summary level, overall and in three document-length buckets.
\item Evidence that both metrics keep the human ranking of systems at every length, while the summary-level discrimination of SummaC weakens on the longest documents and that of QAFactEval does not.
\end{itemize}
